\documentclass[11pt]{article}
\usepackage{mathblatt}
\begin{document}
\blattkopf*{Basisaufgaben}{BAS-Z4}{Basisaufgaben · Zettel · BAS-Z4}
\einheitenkopf[][Zettel 4]{Basisaufgaben · Zettel 4}
\zweigzeile{ohne Taschenrechner · je Aufgabe 1 Punkt · Lösungen auf der Rückseite}
% \small: zehn Aufgaben mit bis zu zwei Koordinatensystemen auf einer Seite (Render 28.09.)
\small

% 1: Mitte zweier Zahlen bestimmen – brueche-dezimalzahlen-basis-k3-v1 (Original 2020-OS-B1f)
\begin{aufgabe}{Welche Zahl liegt genau in der Mitte zwischen $-0{,}4$ und $-0{,}3$? \hfill (P20) \\ \leerfeld}
\end{aufgabe}

% 2: Prozentwert berechnen – prozentrechnung-basis-k4-v3 (Original 2026-FOR-B1a)
\begin{aufgabe}{$25\,\%$ von $64\,€$ \hfill (P26F) \\ \leerfeld[€]}
\end{aufgabe}

% 3: Term durch Zusammenfassen gleichartiger Glieder vereinfachen – terme-basis-k1-v1 (Original 2025-GYM-B2a)
\begin{aufgabe}{Vereinfache den Term $7x - 2x^2 - 3x$ und berechne seinen Wert für $x = 1$. \hfill (P25G)}
\end{aufgabe}

% 4: Term zu Sachtext angeben – terme-basis-k2-v2 (Original 2023-OS-B1h)
\begin{aufgabe}{Die Differenz aus dem Fünffachen einer Zahl $a$ und 6 wird halbiert. Welcher Term passt? \hfill (P23) \\ \kreuz{$2 \cdot (5a - 6)$} \kreuz{$5a - 6 : 2$} \kreuz{$(5a - 6) : 2$} \kreuz{$5a : 6 - 2$}}
\end{aufgabe}

% 5: Lösung durch Einsetzen prüfen – quadratische-gleichungen-basis-k1-v2 (Original 2025-OS-B1h)
\begin{aufgabe}{Welcher Wert erfüllt $x \cdot (x + 4) = -3$? \hfill (P25) \\ \kreuz{$x = 1$} \kreuz{$x = 3$} \kreuz{$x = -3$} \kreuz{$x = -4$}}
\end{aufgabe}

% 6: Graph nach Eigenschaft auswählen – lineare-funktionen-basis-k2-v1 (Original 2016-OS-B1c)
\begin{aufgabe}{\begin{minipage}[t]{0.6\linewidth}Welche Gerade steigt? Kreuze an. \hfill (P16) \\ \kreuz{f} \kreuz{g}\end{minipage}\hfill\begin{minipage}[t]{0.37\linewidth}\vspace{0pt}
\begin{ksys}[xmin=-5,xmax=5,ymin=-5,ymax=5,klein] \gerade{-2}{1}{f} \gerade{0.5}{-1}{g} \end{ksys}
\end{minipage}}
\end{aufgabe}

% 7: Winkel an geschnittenen Parallelen bestimmen – winkel-dreiecke-basis-k3-v1 (Original 2015-OS-B1d)
\begin{aufgabe}{\begin{minipage}[t]{0.6\linewidth}Die beiden waagerechten Geraden sind parallel. Wie groß ist $\alpha$? \hfill (P15) \\ $\alpha =$ \leerfeld °\end{minipage}\hfill\begin{minipage}[t]{0.37\linewidth}\vspace{0pt}
\parallelenpaar{62}{62^\circ}{}{}{\alpha}
\end{minipage}}
\end{aufgabe}

% 8: Winkel im Viereck berechnen – winkel-dreiecke-basis-k4-v1 (Original 2026-FOR-B1i)
\begin{aufgabe}{\begin{minipage}[t]{0.6\linewidth}Ein Parallelogramm hat links unten einen Winkel von $68^\circ$. Wie groß ist der Winkel $\beta$ rechts unten? \hfill (P26F) \\ $\beta =$ \leerfeld °\end{minipage}\hfill\begin{minipage}[t]{0.37\linewidth}\vspace{0pt}
\parallelogramm[winkel={68^\circ,\beta,,}]{3}{1.8}{68}
\end{minipage}}
\end{aufgabe}

% 9: Trigonometrische Gleichung nach Seite umstellen – trigonometrie-basis-k1-v1 (Original 2020-OS-B1j)
\begin{aufgabe}{$\mathrm{sin}\,30^\circ = \frac{5}{x}$ – berechne $x$. \hfill (P20) \\ x = \leerfeld}
\end{aufgabe}

% 10: Winkelfunktion Seitenverhältnis angeben – trigonometrie-basis-k2-v2 (Original 2025-OS-B1g)
\begin{aufgabe}{\begin{minipage}[t]{0.6\linewidth}Rechter Winkel bei $A$. Trage den Bruch ein. \hfill (P25) \\ $\mathrm{cos}\,\gamma =$ \leerfeld\end{minipage}\hfill\begin{minipage}[t]{0.37\linewidth}\vspace{0pt}
\dreieck{(0,0)}{(3,0)}{(0,2.2)}{p}{q}{r}{}{}{\gamma}
\end{minipage}}
\end{aufgabe}

\begleitteil
\erg{1}{$(-0{,}4 + (-0{,}3)) : 2 = -0{,}35$}
\erg{2}{$0{,}25 \cdot 64 = 16\,€$}
\erg{3}{$7x - 3x = 4x$, Term $4x - 2x^2$; für $x = 1$: $4 \cdot 1 - 2 \cdot 1^2 = 4 - 2 = 2$}
\erg{4}{$(5a - 6) : 2$}
\erg{5}{$x = -3$, denn $(-3) \cdot (-3 + 4) = (-3) \cdot 1 = -3$}
\erg{6}{g (von links nach rechts gelesen geht sie nach oben)}
\erg{7}{$\alpha = 180^\circ - 62^\circ = 118^\circ$ (Stufenwinkel, dann Nebenwinkel)}
\erg{8}{$\beta = 180^\circ - 68^\circ = 112^\circ$}
\erg{9}{$x = 5 : \mathrm{sin}\,30^\circ = 5 : 0{,}5 = 10$}
\erg{10}{$\mathrm{cos}\,\gamma = \frac{q}{p}$}
\end{document}
